%% LyX 2.4.4 created this file.  For more info, see https://www.lyx.org/.
%% Do not edit unless you really know what you are doing.
\documentclass[english]{article}
\usepackage[T1]{fontenc}
\usepackage[latin9]{inputenc}
\usepackage{xcolor}
\usepackage{booktabs}
\usepackage{calc}
\usepackage{units}
\usepackage{amsmath}
\usepackage{amsthm}
\usepackage{cancel}
\usepackage{geometry}
\geometry{verbose,tmargin=1in,bmargin=1in,lmargin=1in,rmargin=1in}

\makeatletter

%%%%%%%%%%%%%%%%%%%%%%%%%%%%%% LyX specific LaTeX commands.
%% Because html converters don't know tabularnewline
\providecommand{\tabularnewline}{\\}

%%%%%%%%%%%%%%%%%%%%%%%%%%%%%% Textclass specific LaTeX commands.
\numberwithin{equation}{section}
\numberwithin{figure}{section}
\newlength{\lyxlabelwidth}      % auxiliary length 

%%%%%%%%%%%%%%%%%%%%%%%%%%%%%% User specified LaTeX commands.
\usepackage{fancyhdr}
\usepackage{lastpage}
\pagestyle{fancy}
\usepackage{enumitem}
\usepackage{ifthen}
%sets math fonts to be more readable
\everymath=\expandafter{\the\everymath\displaystyle}
%\DeclareMathSizes{10}{10}{10}{10}
\usepackage{multicol}

\usepackage{tikz}
\usetikzlibrary{plotmarks}
\usepackage{pgfplots}
\pgfplotsset{compat=newest}

\usepackage{calc}
\usepackage[nomessages]{fp}

\usepackage{advdate}    % Advancing/saving dates
\usepackage{datetime}   % Dates formatting
\usepackage{datenumber} % Counters for dates
% Added by lyx2lyx
\setlength{\parskip}{\smallskipamount}
\setlength{\parindent}{0pt}

\makeatother

\usepackage{babel}
\begin{document}
\renewcommand{\author}{Dr.\,James Ashe}

\newcommand{\classNum}{132}

\newcommand{\classSec}{C and K}

\newcommand{\nameType}{}

\newcommand{\examType}{Formula Sheet}

\renewcommand{\date}{}

%%First page's header%%%%%%%%%%%%%%%%%%%%%%
\fancypagestyle{firststyle} %%header settings for first pagr
{
	\fancyhead[L]{MTH \classNum{}(\classSec{}) \author{}\\
	\textbf{\examType{}} \date{}}
	\fancyhead[C]{\textbf{\nameType}}
	\setlength{\headsep}{14pt}
}
%%%%%%%%%%%%%%%%%%%%%%%%%%%%%%%%%%%%%%%%%%

%%Next page's header%%%%%%%%%%%%%%%%%%%%%%%
\setlist{nolistsep}
\lhead{MTH \classNum{}(\classSec{})} 
\chead{\textbf{\nameType}} 
\cfoot{\thepage{}~of~\pageref{LastPage}} 
\setlength{\headsep}{5pt} 
%%%%%%%%%%%%%%%%%%%%%%%%%%%%%%%%%%%%%%%%%%%

%%Various settings; see the preamble for other settings%%%%%
%%\setenumerate[0]{leftmargin=12pt,itemindent=0pt,labelwidth=0pt}
\renewcommand{\labelenumi}{\textbf{\arabic{enumi}.}}
\renewcommand{\labelenumii}{\textbf{\alph{enumii}.}}
\thispagestyle{firststyle}
%%%%%%%%%%%%%%%%%%%%%%%%%%%%%%%%%%%%%%%%%%
\newcommand{\xmax}{10}
\newcommand{\xmin}{-10}
\newcommand{\xstep}{0} %must be xmin+n
\newcommand{\ymax}{10}
\newcommand{\ymin}{-10}
\newcommand{\ystep}{0} %must be ymin+n

\newcommand{\xspace}{\vspace{1.5cm}}

\subsection*{\center{Formula Sheet}}

\noindent\fcolorbox{black}{white}{\begin{minipage}[t]{1\columnwidth - 2\fboxsep - 2\fboxrule}%
\begin{center}
\begin{tabular}{ccccc}
\textbf{Slope-intercept form} &  & \textbf{Slope formula} &  & \textbf{Point-Slope formula}\tabularnewline
$y=mx+b$ &  & $m=\frac{y_{2}-y_{1}}{x_{2}-x_{1}}$ &  & $y-y_{1}=m(x-x_{1})$\tabularnewline
\end{tabular}
\par\end{center}%
\end{minipage}}

\vfill%
\noindent\begin{minipage}[t]{1\columnwidth}%
\begin{center}
\begin{tabular}{llll}
$P=$principal/present value & $r=$annual interest rate & $m=$compoundings per year & $n=tm$\tabularnewline
$A=$future/maturity value & $t=$number of years & $i=\nicefrac{r}{m}=$rate per period & $r_{E}=$effective rate\tabularnewline
\end{tabular}
\par\end{center}%
\end{minipage}

\noindent\fcolorbox{black}{white}{\begin{minipage}[t]{1\columnwidth - 2\fboxsep - 2\fboxrule}%
\begin{center}
\begin{tabular}{lclcl}
\textbf{Simple Interest} &  & \textbf{Compound Interest} &  & \textbf{Continuous Compounding}\tabularnewline
\midrule
$A=P(1+rt)$ &  & $A=P(1+i)^{n}$ &  & $A=Pe^{rt}$\tabularnewline
$P=\frac{A}{1+rt}$ &  & $P=A(1+i)^{-n}$ &  & $P=Ae^{-rt}$\tabularnewline
 &  & $r_{E}=\left(1+\frac{r}{m}\right)^{m}-1$ &  & $r_{E}=e^{r}-1$\tabularnewline
\end{tabular}
\par\end{center}%
\end{minipage}}

\vfill%
\noindent\begin{minipage}[t]{1\columnwidth}%
\begin{center}
\begin{tabular}{lll}
$R=$periodic payment & $P=$present value of annuity & $n=$number of periods\tabularnewline
$S=$future value &  & $i=$rate per period\tabularnewline
\end{tabular}
\par\end{center}%
\end{minipage}

\noindent\fcolorbox{black}{white}{\begin{minipage}[t]{1\columnwidth - 2\fboxsep - 2\fboxrule}%
\begin{center}
\begin{tabular}{lcl}
\textbf{Future Value of an Annuity} &  & \textbf{Sinking Fund Payment}\tabularnewline
\midrule
$S=R\left[\frac{\left(1+i\right)^{n}-1}{i}\right]$ &  & $R=\frac{S\cdot i}{(1+i)^{n}-1}$\tabularnewline
\end{tabular}
\par\end{center}%
\end{minipage}}

\noindent\fcolorbox{black}{white}{\begin{minipage}[t]{1\columnwidth - 2\fboxsep - 2\fboxrule}%
\begin{center}
\begin{tabular}{ll}
\textbf{Present Value of an Annuity} & \textbf{Amortization Payments}\tabularnewline
\midrule
$P=R\left[\frac{1-\left(1+i\right)^{-n}}{i}\right]$ & $R=\frac{P}{\left[1-(1+i)^{-n}\right]}$\tabularnewline
\end{tabular}
\par\end{center}%
\end{minipage}}

\vfill

\begin{center}
\textbf{Ti-83/84 }
\par\end{center}

\begin{center}
\begin{tabular}{rl}
\textbf{N} & : number of payments\tabularnewline
${\tt I\%}$ & : interest rate\tabularnewline
${\tt PV}$ & : present value\tabularnewline
${\tt PMT}$ & : payment size\tabularnewline
${\tt FV}$ & : future value\tabularnewline
${\tt P/Y}$ & : payments per year\tabularnewline
${\tt C/Y}$ & : compoundings per year\tabularnewline
\end{tabular}
\par\end{center}

\begin{center}
\vfill
\par\end{center}
\end{document}
